
\subsection{算子 $\mathfrak q$ 的构造 III：
$\operatorname{Der}(\widehat B(C[1]),\widehat B(C[1]))$
的同伦不变性}
\label{subsec:7-4-4}

设 $(C,\mathfrak m)$ 与
$(C',\mathfrak m')$ 为有隙滤过
$A_\infty$ 代数。
把微分分次李代数
$\operatorname{Der}B(C[1])$、
$\operatorname{Der}B(C'[1])$
视为滤过 $L_\infty$ 代数。

\begin{theorem}{7.4.41}\label{thm:7-4-41}
若 $(C,\mathfrak m)$ 与
$(C',\mathfrak m')$
作为滤过 $A_\infty$ 代数
有隙同伦等价，
则
$\operatorname{Der}B(C[1])$
与
$\operatorname{Der}B(C'[1])$
作为滤过 $L_\infty$ 代数
严格且有隙地同伦等价。
含幺版本亦成立。
\end{theorem}

\begin{remark}{7.4.42}\label{rem:7-4-42}
\par\noindent\textup{(1)}\quad
至少对非滤过版本而言，
这一结果似为专家所知。
事实上，定理~\ref{thm:7-4-41}
可视为对既知的形变函子同伦等价
给出一种系统表述。

\par\noindent\textup{(2)}\quad
若进一步在 $C$ 上有内积，
且 $\mathfrak m$ 满足循环对称性
（见第~\ref{subsec:3-6-4} 小节），
则可在
\[
\operatorname{Hom}^{\mathrm{filt}}
\bigl(B^{\mathrm{cyc}}(C[1]),
\Lambda_{0,\mathrm{nov}}\bigr)
\]
上定义微分分次李代数结构。
由此也可得到类似
定理~\ref{thm:7-4-41} 的结果。
事实上，至少在
$\mathfrak m_k=0$（$k\ne1,2$）
的情形下，
还可证明其“高亏格类比”。
这将联系到本书论述的高亏格推广、
弦拓扑 \textbf{[ChSu99]}，
以及微扰陈--西蒙斯规范场论
\textbf{[AxSi91I,II]}。
这些联系留待别处讨论
（另见注 7.4.97）。

\par\noindent\textup{(3)}\quad
似乎不存在自然的微分分次李代数同态
连接
$\operatorname{Der}B(C[1])$
与
$\operatorname{Der}B(C'[1])$
并诱导同伦等价。
因此，若要证明类似
定理~\ref{thm:7-4-41} 的结果，
似乎不可避免地要使用
$L_\infty$ 代数的语言。
\end{remark}

\par\noindent\textit{定理~\ref{thm:7-4-41} 的证明.}\quad
若 $C$ 与 $C'$ 都是典范模型，
则由命题~\ref{prop:5-4-5}
存在同构
\[
\widehat\varphi:
\widehat B(C[1])
\longrightarrow
\widehat B(C'[1]),
\]
它既是上链映射，
又是余代数同态。
因而
$\operatorname{Der}B(C[1])$
与
$\operatorname{Der}B(C'[1])$
作为微分分次李代数同构，
定理在此情形下成立。

所以只需证明
$C'=H(\overline C,\overline{\mathfrak m}_1)
\otimes_{\mathbb Q}
\Lambda_{0,\mathrm{nov}}$
是 $C$ 的典范模型的情形
（这里
$\overline C\otimes_{\mathbb Q}
\Lambda_{0,\mathrm{nov}}=C$）。
令
$\overline H
=H(\overline C,\overline{\mathfrak m}_1)$、
$H=\overline H\otimes_{\mathbb Q}
\Lambda_{0,\mathrm{nov}}$，
并如定理~\ref{thm:5-4-2}
的证明那样在 $H$ 上定义
滤过 $A_\infty$ 结构。
我们将构造滤过 $L_\infty$ 同态
\begin{equation}
\widehat{\mathfrak g}^{\circ}:
\operatorname{Der}B(C[1])
\longrightarrow
\operatorname{Der}B(H[1]).
\tag{7.4.43}\label{eq:7-4-43}
\end{equation}
按定义，
$\mathfrak g^{\circ}$
可等同于一族次数为 $1$
的滤过同态
\begin{equation}
\begin{aligned}
\mathfrak g^{\circ}_{\ell,k}:
E_\ell\bigl(
\operatorname{Der}B(C[1])[1]
\bigr)
\otimes
\widehat B_k(H[1])
\longrightarrow H[1].
\end{aligned}
\tag{7.4.44}\label{eq:7-4-44}
\end{equation}
进一步分解为
\[
\mathfrak g^{\circ}_{\ell,k}
=\sum_{\beta\in G}
T^{E(\beta)}e^{\mu(\beta)/2}
\mathfrak g^{\circ}_{\ell,k,\beta},
\]
其中
\[
G=\bigl\{
(E(\beta),\mu(\beta))
\in\mathbb R_{\ge0}\times2\mathbb Z
\mid\beta\in G
\bigr\}
\]
是满足条件 3.1.6 的
$\mathbb R_{\ge0}\times2\mathbb Z$
的子幺半群，并且
\begin{equation}
\begin{aligned}
\mathfrak g^{\circ}_{\ell,k,\beta}:
E_\ell\bigl(
\operatorname{Der}B(\overline C[1])[1]
\bigr)
\otimes
B_k(\overline H[1])
\longrightarrow\overline H[1]
\end{aligned}
\tag{7.4.45}\label{eq:7-4-45}
\end{equation}
是次数为 $1$ 的
$\mathbb Q$-线性映射。
仍如前文，假定 $C$ 与 $H$
均为 $G$-有隙。

式~\eqref{eq:7-4-46}
的定义将由对树求和给出，
与第~\ref{subsec:5-4-4} 小节相似。\translatornote{源书在此称“式 (7.4.46) 的定义”由树求和给出；但该式实际上直接定义置换对称化 $\operatorname{Sym}$，树的递归构造从后文才开始。这里保留原指称，见疑误 A-0565。}
设
\[
Y_a\in
\operatorname{Hom}
\bigl(B_{k_a}(\overline C[1]),
\overline C[1]\bigr)
\subset
\operatorname{Hom}
\bigl(B(\overline C[1]),
\overline C[1]\bigr)
\cong
\operatorname{Der}B(\overline C[1]).
\]
记
\begin{equation}
\operatorname{Sym}(Y_1,\ldots,Y_\ell)
=\frac1{\ell!}
\sum_{\sigma\in\mathfrak S_\ell}
(-1)^{\ast(\sigma)}
Y_{\sigma(1)}\otimes\cdots\otimes
Y_{\sigma(\ell)},
\tag{7.4.46}\label{eq:7-4-46}
\end{equation}
其中指数约定为\translatornote{源书以输出次序
$Y_{\sigma(1)}\otimes\cdots\otimes
Y_{\sigma(\ell)}$
配用下式指数；在混合奇偶次数下，
它与按该输出顺序直接计算的
Koszul 符号可能不同。
置换作用约定待核，译稿暂保留源式，见疑误 A-0566。}：
\[
\ast(\sigma)
=\sum_{\substack{i<j\\\sigma(i)>\sigma(j)}}
\deg'Y_i\,\deg'Y_j.
\]
